% Complete Chapter 5 introduction, source PDF 61 / printed 51.
% block: approximation-hilbert-heading page:61 type:chapter
\clearpage\setcounter{section}{4}
\section{{Approximation of Hilbert complexes}}\label{chap:approximation-hilbert-complexes}

% block: approximation-hilbert-goal page:61 type:ordinary-paragraph
{Our goal now is to design and analyze numerical methods for the approximate solution of a Hodge Laplacian problem associated to some closed Hilbert complex. We shall use Galerkin methods, i.e., we shall replace the trial and test spaces occuring in the primal or mixed weak formulation by finite dimensional subspaces. In practice, when the Hilbert complex is the de Rham complex or something similar, the subspaces we consider will be finite element spaces.}

% block: approximation-short-complex page:61 type:ordinary-paragraph
{To define the Hodge Laplace problem at some level $k$ of the Hilbert complex, we require only the operators $d^{k-1}$ from $W^{k-1}$ to $W^k$ with domain $V^{k-1}$ and the $d^k$ from $W^k$ to $W^{k+1}$ with domain $V^k$; cf. \eqref{eq:abstract-hodge-operator}. Therefore, in the interest of carrying no more baggage than necessary, in this chapter we do not require the whole Hilbert complex \eqref{eq:hilbert-domain-complex} but only a short segment of it:}
\begin{equation}\label{eq:approximation-short-complex}
V^{k-1}\xrightarrow{d}V^k\xrightarrow{d}V^{k+1}.
\end{equation}

% block: approximation-chapter-outline page:61 type:ordinary-paragraph
{We shall see that Galerkin discretization of the primal formulation is problematic and shall focus on Galerkin discretization of the mixed formulation. This requires two compatible Galerkin subspaces, one for $V^{k-1}$ and one for $V^k$. The main result of this chapter is to establish assumptions on the choice of subspaces, which ensures consistency and stability of the Galerkin method, and to deduce consequences from it. Box~\ref{box:approximation-hilbert-summary} summarizes the main results of the chapter.}
